
\subsection{Per Time Calculation (PETIC)}

O Per Time Calculation (PETIC) é um algoritmo proposto para realizar a alocação de fatias considerando as variações na demanda de largura de banda dos clientes ao longo do tempo. Para cada período analisado, o algoritmo busca definir uma fatia capaz de atender aos requisitos do cliente, selecionando os recursos disponíveis na infraestrutura física. Dessa forma, seu processo de alocação considera o estado dos recursos da rede e a demanda apresentada em cada período \cite{gomes2018fatiamento}.

O PETIC é caracterizado como um algoritmo guloso e determinístico, realizando suas decisões de alocação a partir das condições apresentadas pela infraestrutura em cada momento. Para definir uma fatia, o algoritmo considera a largura de banda disponível nos enlaces e a largura de banda requisitada pelo cliente, atribuindo pesos aos recursos que podem atender à solicitação. A partir desses pesos, é realizado um processo de busca para encontrar os caminhos que formarão a fatia \cite{gomes2018fatiamento}.

Durante o processo de alocação, os enlaces que não possuem largura de banda suficiente para atender à requisição são desconsiderados. Para os enlaces restantes, o cálculo dos pesos considera a relação entre a largura de banda disponível e a largura de banda requisitada, favorecendo a utilização de recursos que apresentem maior capacidade disponível. Com isso, o algoritmo busca evitar a saturação dos enlaces e preservar recursos que possam ser utilizados em novas solicitações \cite{gomes2018fatiamento}.

Além da disponibilidade dos recursos no momento da alocação, o PETIC considera informações da fatia utilizada anteriormente. Caso a demanda do período corrente não supere a largura de banda alocada à fatia no período anterior, os enlaces dessa fatia recebem um peso muito próximo de zero, favorecendo sua reutilização. Essa característica busca diminuir alterações na estrutura da fatia entre diferentes períodos, evitando modificações desnecessárias quando a configuração anterior continua adequada às necessidades do cliente \cite{gomes2018fatiamento}.

Do ponto de vista operacional, o PETIC executa, para cada período, os seguintes passos: (i) atribui um peso a cada enlace, conforme a relação entre a banda disponível e a banda requisitada, ou o peso reduzido quando o enlace pertence à fatia anterior reaproveitável; (ii) descarta os enlaces cuja banda disponível é insuficiente para a demanda; e (iii) executa uma busca de menor caminho, no estilo do algoritmo de Dijkstra, sobre os enlaces restantes, retornando o caminho de menor custo total. Se não existir caminho, a requisição é considerada não atendida naquele período. O caminho encontrado é então armazenado e serve de base para a alocação do período seguinte \cite{gomes2018fatiamento}.
